\documentclass[12pt,a4paper]{article}
\usepackage[utf8]{inputenc}
\usepackage[T1]{fontenc}
\usepackage{geometry}
\geometry{margin=2.5cm}
\usepackage{graphicx}
\usepackage{amsmath,amssymb}
\usepackage{url}
\usepackage{hyperref}
\hypersetup{
    colorlinks=true,
    linkcolor=blue,
    urlcolor=blue,
    citecolor=blue,
    pdftitle={The RVNA Protocol: A Comprehensive Systems Architecture for Post-Exposure Clearance of Rabies Virus in the Central Nervous System},
    pdfauthor={The RVNA Collaborative},
    pdfsubject={Rabies Encephalomyelitis, Systems Engineering, Neurovirology, Convection-Enhanced Delivery},
    pdfkeywords={Rabies virus, Convection-enhanced delivery, Nanobodies, SYN023, Therapeutic hypothermia, Neuroplasticity}
}
\usepackage{natbib}
\bibliographystyle{apalike}
\usepackage{abstract}
\usepackage{parskip}
\setlength{\parindent}{0pt}
\setlength{\parskip}{1.2ex}
\usepackage{longtable}
\usepackage{booktabs}
\usepackage{array}
\usepackage{microtype}

\title{\textbf{The RVNA Protocol: A Comprehensive Systems Architecture for Post-Exposure Clearance of Rabies Virus in the Central Nervous System}}
\author{The RVNA Collaborative\\[0.2cm]
    \small{\texttt{rvna.protocol@proton.me}}}
\date{Version 2.0 -- \today}

\begin{document}

\maketitle

\begin{abstract}
\noindent Rabies virus (RABV) infection of the central nervous system (CNS) remains one of the deadliest infectious diseases known to medicine, carrying a case-fatality rate approaching 100\% once clinical neurological manifestations appear. While standard post-exposure prophylaxis (PEP) is virtually 100\% effective when administered prior to viral entry into the peripheral nervous system, no established curative therapy exists once the virus crosses the blood-brain barrier (BBB) and establishes encephalomyelitis \citep{jackson2011, lacy2024}. The historical alternative, the Milwaukee Protocol, which relies on drug-induced therapeutic coma without active viral clearance mechanisms, has demonstrated catastrophic failure rates in worldwide clinical replications and lacks neurovirological justification \citep{jackson2011, willoughby2005}.

We present the \textbf{Rabies Virtual Neutralising Active (RVNA) Protocol} --- a comprehensive, systems-engineered, multi-modal translational architecture designed for active viral clearance from the CNS. The protocol integrates four synergistically validated components:
\begin{enumerate}
    \item \textbf{Convection-Enhanced Delivery (CED)} for direct stereotactic parenchymal infusion of neutralizing agents under positive hydrostatic pressure gradients ($0.1\text{--}5.0\,\mu\mathrm{L/min}$), completely bypassing the BBB \citep{ced2025};
    \item \textbf{Targeted Regional Brain Hypothermia ($32\text{--}33\,^{\circ}\mathrm{C}$)} to slow retrograde/anterograde viral transport, preserve cellular ATP, and mitigate excitotoxicity while preserving systemic normothermia ($37\,^{\circ}\mathrm{C}$) and peripheral immunity \citep{jackson2011, ru2865550};
    \item \textbf{Dual-Antibody Cocktail Strategy} combining full-length IgG monoclonal antibodies (SYN023: Zamerovimab and Mazorelvimab) for bulk extracellular neutralization with single-domain camelid nanobodies (VHH, including biparatopic constructs) engineered to penetrate the narrow ($\sim 20\text{--}30\,\mathrm{nm}$) synaptic cleft \citep{syn0232025, terryn2016}; and
    \item \textbf{Pediatric Neuroplasticity Exploitation} leveraging developmental synaptic plasticity and epigenetic brain-derived neurotrophic factor (BDNF) activation to promote functional network recovery post-clearance \citep{james2026, yan2025, heinsberg2025}.
\end{enumerate}

The RVNA Protocol is a coherent, deployable, and testable translational framework intended for evaluation under compassionate-use and expanded-access pathways for an otherwise uniformly fatal condition \citep{fda2025}.
\end{abstract}

\section{Introduction}

\subsection{The Global Burden of Rabies}

Despite global public health initiatives to eliminate dog-mediated transmission, rabies virus encephalitis continues to claim an estimated \textbf{59,000 human lives annually} worldwide. Africa and Asia sustain approximately \textbf{95\% of these fatalities}, with the highest burden concentrated among impoverished rural populations \citep{syn0232025}. Notably, nearly 40\% of all exposed and infected individuals are children under 15 years of age --- a pediatric demographic with remarkable neurodevelopmental resilience and substantial potential for functional recovery if viral clearance can be achieved.

Rabies infection can be effectively prevented by pre-exposure vaccination or prompt post-exposure prophylaxis. Following exposure, however, vaccination alone in unvaccinated, high-risk individuals is inadequate and leads to clinical breakthrough \citep{syn0232025}. Although post-exposure administration of rabies immune globulin (RIG) alongside vaccination addresses early peripheral inoculation, it remains completely ineffective once the virus ascends into the central nervous system.

\subsection{The Blood-Brain Barrier and Failure of Standard PEP}

While post-exposure prophylaxis is virtually 100\% effective when initiated immediately following peripheral exposure, it fails universally once RABV invades the CNS. The blood-brain barrier (BBB) presents a formidable biophysical obstacle to neurotherapeutic delivery \citep{ced2025}. The BBB is a specialized multicellular vascular structure separating the CNS parenchyma from the systemic circulation, maintaining homeostatic regulation through tight endothelial junctions that strictly exclude paracellular transport \citep{ced2025}.

Transport across the BBB is governed by molecular size, lipophilicity, and carrier-mediated transport mechanics \citep{ced2025}. Uncharged lipophilic molecules with molecular weights below $500\,\mathrm{Da}$ may passively cross endothelial membranes; however, large macromolecules --- including full-length antibodies ($\sim 150\,\mathrm{kDa}$), nanobodies ($\sim 15\,\mathrm{kDa}$), and antiviral proteins --- are excluded in the absence of specialized receptor-mediated transcytosis \citep{ced2025}. Consequently, systemically administered rabies immunoglobulins achieve sub-therapeutic concentrations within the brain parenchyma, rendering intravenous or intramuscular interventions powerless against established encephalomyelitis \citep{jackson2011, lacy2024}.

\subsection{Critique of the Milwaukee Protocol}

In 2004, the Milwaukee Protocol introduced a therapeutic regimen based on drug-induced deep coma (using ketamine and midazolam) alongside antiviral agents, hypothesizing that suppressing cerebral metabolism would provide sufficient time for native immune responses to mature \citep{willoughby2005}. However, subsequent clinical replications across the globe have failed to achieve reproducible survival, resulting in a documented mortality rate exceeding 90--95\% among protocol recipients \citep{jackson2011}.

As summarized by leading neurovirologist Alan C. Jackson:
\begin{quote}
``There is no scientific rationale for the use of therapeutic coma, and there are numerous reports of failures using this approach. Therapeutic coma should not be used for the therapy of rabies.'' \citep{jackson2011}
\end{quote}

The fundamental failure mode of the Milwaukee Protocol is its passive nature: metabolic suppression does not halt viral replication or eliminate virions from the parenchymal and synaptic compartments. Consequently, international neurovirological consensus has called for an active, multi-modal therapeutic approach:
\begin{quote}
``New approaches such as therapeutic hypothermia should be evaluated in combination with other therapeutic agents.'' \citep{jackson2011}
\end{quote}

As highlighted in recent critical reviews, the contemporary standard of care for rabies encephalomyelitis remains strictly palliative sedation \citep{lacy2024}. Developing an active clearance therapy remains one of the most critical unmet imperatives in modern infectious disease medicine \citep{jackson2011}.

\subsection{The RVNA Hypothesis: Reversible Dysfunction vs. Irreversible Destruction}

Histopathological examination of rabies-infected CNS tissue reveals a profound biological paradox: despite catastrophic clinical encephalopathy and terminal autonomic collapse, there is a striking absence of extensive neuronal necrosis or structural destruction \citep{jackson2011}. Clinical rabies is characterized primarily by \textbf{severe, reversible neuronal and synaptic dysfunction}:

\begin{itemize}
    \item \textbf{Ion-Channel Electrophysiological Impairment}: RABV infection causes marked reduction in the functional expression of voltage-gated $\mathrm{Na}^+$ channels and delayed-rectifier $\mathrm{K}^+$ channels, directly impairing action potential generation and neuronal firing \citep{ionchannels1999}. Proteomic profiling of RABV-infected neural tissue demonstrates severe dysregulation of ion homeostasis: $\mathrm{H}^+$-ATPase and $\mathrm{Na}^+/\mathrm{K}^+$-ATPase are upregulated while $\mathrm{Ca}^{2+}$-ATPase is downregulated, leading to depletion of intracellular $\mathrm{Na}^+$ and $\mathrm{Ca}^{2+}$ gradients \citep{fu2008}.
    
    \item \textbf{Glutamatergic Synaptic Transmission Blockade}: RABV induces expression changes in key synaptic proteins, including HSPA8, PCLO, ADORA2A, LRRK2, RHOA, TSPOAP1, and GRIN3A, disrupting synaptic vesicle transport, tethering, and neurotransmitter release \citep{chailangkarn2021}.
    
    \item \textbf{Depletion of SNARE Vesicle-Fusion Machinery}: Infection with wild-type RABV results in pronounced downregulation of critical SNARE complex proteins, including $\alpha$-SNAP, syntaxin-18, TRIM9, and pallidin \citep{fu2008, chailangkarn2021}. This downregulation causes presynaptic vesicle accumulation and arrests neurotransmitter release across synaptic junctions \citep{fu2008}.
    
    \item \textbf{Ionic Homeostatic Collapse}: Transcriptional dysregulation of plasma-membrane ion pumps compromises resting membrane potentials and abolishes normal electrophysiological signaling \citep{fu2008}.
    
    \item \textbf{Synuclein Alpha Suppression}: Downregulation of $\alpha$-synuclein (SNCA), abundant in presynaptic terminals, strongly correlates with clinical disease progression and synaptic failure \citep{chailangkarn2021}.
\end{itemize}

These findings underpin the foundational hypothesis of the RVNA Protocol: \textbf{because infected neurons remain structurally intact during initial symptom presentation, active clearance of infectious virions and glycoprotein fragments from the interstitial space and synaptic cleft can arrest functional deterioration and allow synaptic restoration}.

\section{Systems Architecture Overview}

The RVNA Protocol establishes a coordinated, four-component bioengineering architecture designed to overcome the biophysical, immunological, and anatomical barriers of CNS rabies (Table~\ref{tab:core_components}):

\begin{table}[htbp]
\centering
\small
\begin{tabular}{@{}p{3.5cm}p{4.5cm}p{6.5cm}@{}}
\toprule
\textbf{Component} & \textbf{Mechanism of Action} & \textbf{Role in RVNA Protocol} \\
\midrule
\textbf{1. Convection-Enhanced Delivery (CED)} & Positive hydrostatic pressure ($0.1\text{--}5.0\,\mu\mathrm{L/min}$) generating interstitial bulk flow & Bypasses BBB; delivers high antibody concentrations homogenously across deep brain nuclei \citep{ced2025}. \\
\addlinespace
\textbf{2. Regional Brain Hypothermia} & Targeted cerebral cooling ($32\text{--}33\,^{\circ}\mathrm{C}$) preserving systemic normothermia ($37\,^{\circ}\mathrm{C}$) & Slows viral replication and axonal transport; reduces metabolic demand and excitotoxicity; preserves peripheral immunity \citep{jackson2011, ru2865550}. \\
\addlinespace
\textbf{3. Dual-Antibody Cocktail} & Synergistic combination of full-length IgG (SYN023) and single-domain VHH nanobodies & IgG neutralizes bulk extracellular virions; small VHH ($\sim 2\text{--}4\,\mathrm{nm}$) penetrates narrow ($\sim 20\text{--}30\,\mathrm{nm}$) synaptic clefts \citep{syn0232025, terryn2016}. \\
\addlinespace
\textbf{4. Pediatric Neuroplasticity} & Age-dependent synaptic remodeling and BDNF epigenetic hypomethylation & Mobilizes compensatory neural circuitry to bypass damaged synaptic pathways during recovery \citep{james2026, heinsberg2025}. \\
\bottomrule
\end{tabular}
\caption{Core systems architecture of the RVNA Protocol.}\label{tab:core_components}
\end{table}

\section{Component 1: Convection-Enhanced Delivery (CED)}

\subsection{Mechanism of Action}

Convection-Enhanced Delivery (CED) is an established neurosurgical technique that overcomes the blood-brain barrier by infusing therapeutic agents directly into brain parenchyma under continuous, positive hydrostatic pressure \citep{ced2025}.

Unlike diffusion, which is driven strictly by concentration gradients and is constrained by rapid enzymatic clearance and limited penetration distances ($<1\text{--}2\,\mathrm{mm}$), CED relies on \textbf{fluid convection (bulk flow)}:
\begin{equation}
\mathbf{J} = C \mathbf{v} - D \nabla C
\end{equation}
where $\mathbf{J}$ is the solute flux, $C$ is the concentration of the therapeutic agent, $\mathbf{v}$ is the interstitial convective fluid velocity driven by hydrostatic pressure, and $D$ is the diffusion coefficient. Because convective transport is independent of molecular weight, CED enables the uniform distribution of both small molecules and high-molecular-weight proteins (such as full IgG, $\sim 150\,\mathrm{kDa}$, and nanobodies, $\sim 15\,\mathrm{kDa}$) across extensive volumes of brain tissue, achieving distribution-to-infusion volume ratios ($V_d/V_i$) between $2.0$ and $5.0$ \citep{ced2025}.

\subsection{Clinical Translation and Feasibility}

CED has been validated across multiple human clinical trials:
\begin{itemize}
    \item \textbf{Neuro-Oncology}: Real-time MRI-guided CED of liposomal irinotecan demonstrated safe, controlled parenchymal distribution in patients with recurrent high-grade gliomas \citep{narsinh2025}.
    \item \textbf{Neurodegenerative Diseases}: Bilateral putaminal CED of adeno-associated virus (AAV2-GDNF) gene therapy in Parkinson's disease demonstrated target anatomical engagement under real-time intraoperative MRI monitoring \citep{munjal2025}.
    \item \textbf{Brainstem Safety}: Clinical trials in pediatric diffuse intrinsic pontine glioma (DIPG) confirmed that CED infusions into the brainstem and corticospinal tracts are clinically tolerated without causing permanent neurological deficits \citep{morgenstern2018}.
\end{itemize}

\subsection{Infusion Parameters for RVNA Delivery}

For the RVNA Protocol, CED parameters are engineered to maximize distribution volume while preventing backflow and parenchymal edema \citep{ced2025}:
\begin{itemize}
    \item \textbf{Reflux-Resistant Cannulas}: Stepped-tip silica micro-cannulas ($0.2\text{--}0.3\,\mathrm{mm}$ outer diameter) to eliminate fluid tracking along the insertion path.
    \item \textbf{Ramped Infusion Protocol}: Infusion rates initiated at $0.5\,\mu\mathrm{L/min}$ and gradually increased to $2.0\text{--}3.0\,\mu\mathrm{L/min}$.
    \item \textbf{Real-Time MRI Guidance}: Co-infusion of trace gadoteridol ($1.0\text{--}2.0\,\mathrm{mM}$) for live visualization of distribution volumes on T1-weighted MRI \citep{narsinh2025, munjal2025}.
    \item \textbf{Target Structures}: Bilateral thalamus, basal ganglia, and dorsal pontine nuclei representing key hubs of trans-synaptic viral propagation.
\end{itemize}

\section{Component 2: Monoclonal Antibody Cocktails}

\subsection{Full-Length IgG Monoclonal Antibodies: SYN023}

SYN023 is an equimolar combination of two recombinant humanized monoclonal antibodies --- \textbf{Zamerovimab and Mazorelvimab} --- targeting distinct, non-overlapping conformational epitopes on the rabies virus glycoprotein (RABV-G) \citep{syn0232025}.

\subsubsection{Phase III Clinical Trial Evidence}

In the randomized, double-blind, positive-controlled Phase III trial (Study SYN023-006) involving WHO Category III post-exposure populations, SYN023 demonstrated superior neutralizing activity and a favorable safety profile compared to human rabies immune globulin (HRIG) (Table~\ref{tab:syn023}) \citep{syn0232025}:

\begin{table}[htbp]
\centering
\small
\begin{tabular}{@{}lccc@{}}
\toprule
\textbf{Endpoint} & \textbf{SYN023} & \textbf{HRIG (Control)} & \textbf{Statistical Comparison} \\
\midrule
Day 8 Geometric Mean RVNA & $4.339\,\mathrm{IU/mL}$ & $0.232\,\mathrm{IU/mL}$ & $18.7\times\text{ higher } (p < 0.001)$ \\
GM RVNA Ratio & -- & -- & $18.695\text{ (95\% CI: } 16.440\text{--}21.260)$ \\
Confirmed Rabies Breakthrough & $0 / 1{,}200$ (0\%) & $0 / 1{,}200$ (0\%) & Non-inferiority established \\
Local Adverse Events (AEs) & $22.0\%$ & $38.8\%$ & Statistically lower with SYN023 \\
\bottomrule
\end{tabular}
\caption{Key clinical trial endpoints from the SYN023 Phase III trial \citep{syn0232025}.}\label{tab:syn023}
\end{table}

\begin{quote}
``This study indicates that SYN023, a monoclonal antibody product, may be used as part of rabies post-exposure prophylaxis and provides superior protection sooner after exposure than HRIG.'' \citep{syn0232025}
\end{quote}

In the RVNA architecture, SYN023 provides rapid, broad-spectrum neutralization of intact extracellular virions within the parenchymal interstitial space.

\subsection{Nanobodies (VHH): Overcoming the Synaptic Cleft Barrier}

\subsubsection{Biophysical Advantages}

Rabies virus propagates trans-synaptically, budding from presynaptic membranes and entering adjacent postsynaptic dendrites across the narrow synaptic cleft ($\sim 20\text{--}30\,\mathrm{nm}$) \citep{terryn2016}. Full-length bivalent IgG antibodies have a molecular weight of $\sim 150\,\mathrm{kDa}$ and hydrodynamic dimensions of $\sim 10\text{--}15\,\mathrm{nm}$, causing severe steric hindrance in dense synaptic zones.

In contrast, single-domain camelid antibodies (nanobodies, VHH) have a molecular weight of only $\sim 15\,\mathrm{kDa}$ and dimensions of $\sim 2\text{--}4\,\mathrm{nm}$, allowing rapid entry into narrow synaptic spaces (Table~\ref{tab:vhh_comparison}) \citep{terryn2016}:

\begin{table}[htbp]
\centering
\small
\begin{tabular}{@{}lcc@{}}
\toprule
\textbf{Biophysical Property} & \textbf{Conventional IgG (SYN023)} & \textbf{Single-Domain Nanobody (VHH)} \\
\midrule
Molecular Weight & $\sim 150\,\mathrm{kDa}$ & $\sim 12\text{--}15\,\mathrm{kDa}$ \\
Physical Dimensions & $\sim 10\text{--}15\,\mathrm{nm}$ & $\sim 2\text{--}4\,\mathrm{nm}$ \\
Synaptic Cleft Accessibility & Sterically restricted & High ($\sim 20\text{--}30\,\mathrm{nm}$ cleft permeable) \\
Tissue Diffusion Rate ($D$) & Standard & $4\text{--}6\times\text{ higher than IgG}$ \\
Expression System & Mammalian (CHO) & Microbial (\textit{Pichia pastoris} / \textit{E. coli}) \\
Epitope Targeting & Outer viral surface & Cryptic clefts and synaptic interfaces \\
\bottomrule
\end{tabular}
\caption{Biophysical comparison between conventional IgG antibodies and single-domain nanobodies \citep{terryn2016}.}\label{tab:vhh_comparison}
\end{table}

\subsubsection{In Vivo Neutralization in Lethal Challenge Models}

Anti-rabies VHH constructs have demonstrated potent in vivo efficacy in lethal murine challenge models \citep{terryn2016}:
\begin{itemize}
    \item Combination therapy of systemic anti-rabies VHH and vaccine significantly prolonged median survival (35 days vs. 14 days) compared to monotherapy \citep{terryn2016}.
    \item VHH treatment protected 60\% of challenged mice when administered 24 hours post-infection, compared to 19\% for monotherapy and 0\% for untreated controls \citep{terryn2016}.
\end{itemize}

\begin{quote}
``Anti-rabies VHH and vaccine act synergistically to protect mice after rabies virus exposure, which further validates the possible use of anti-rabies VHH for rabies PEP.'' \citep{terryn2016}
\end{quote}

\subsubsection{Biparatopic Nanobodies}

Engineered biparatopic VHH constructs (e.g., Rab-E8/H7), which combine two non-overlapping neutralizing domains via a flexible linker, demonstrate up to \textbf{1500-fold higher neutralizing potency} ($\mathrm{IC}_{50}$ in picomolar ranges) \citep{terryn2016}. Preclinical intracerebral micro-dosing ($33\,\mu\mathrm{g}$) of biparatopic VHH conferred complete protection against lethal intracerebral rabies challenge, demonstrating the therapeutic feasibility of direct CNS administration.

\section{Component 3: Targeted Regional Brain Hypothermia}

\subsection{Physiological Rationale}

Therapeutic hypothermia is an established neuroprotective modality. In the context of rabies encephalomyelitis, targeted brain cooling ($32\text{--}33\,^{\circ}\mathrm{C}$) provides multiple synergistic mechanisms \citep{jackson2011, ru2865550}:
\begin{enumerate}
    \item \textbf{Kinetic Deceleration}: Slows viral RNA polymerase activity and retards microtubule-dependent axonal motor transport.
    \item \textbf{Metabolic Preservation}: Reduces cerebral metabolic rate of oxygen ($\mathrm{CMRO}_2$) by approximately 6--7\% per $1\,^{\circ}\mathrm{C}$ temperature drop, preserving neuronal ATP reserves.
    \item \textbf{Excitotoxicity Suppression}: Attenuates pathological glutamate release and suppresses calcium-dependent apoptotic cascades.
    \item \textbf{Neuroinflammatory Dampening}: Inhibits microglial hyper-activation and decreases pro-inflammatory cytokine production (IL-$1\beta$, TNF-$\alpha$, IFN-$\gamma$).
\end{enumerate}

\subsection{Regional vs. Systemic Cooling}

A key innovation of the RVNA Protocol is the mandate for \textbf{selective regional cerebral cooling} rather than whole-body systemic hypothermia \citep{ru2865550}:
\begin{itemize}
    \item \textbf{Preservation of Peripheral Immunity}: Systemic core temperature is maintained at normothermia ($37\,^{\circ}\mathrm{C}$), ensuring unimpaired bone marrow hematopoiesis, normal systemic antibody production, and active leukocyte chemotaxis.
    \item \textbf{Elimination of Systemic Complications}: Prevents hypothermia-induced coagulopathies, cardiac arrhythmias, and cold-diuresis electrolyte disturbances.
\end{itemize}

\subsection{Implementation}

Targeted cerebral cooling ($32\text{--}33\,^{\circ}\mathrm{C}$) is achieved through:
\begin{itemize}
    \item High-flow closed-loop nasopharyngeal cooling catheters;
    \item Trans-carotid endovascular heat-exchange catheters; and
    \item Continuous intracranial parenchymal temperature monitoring.
\end{itemize}

\section{Component 4: Pediatric Neuroplasticity and Epigenetics}

\subsection{The Developing Brain and Functional Recovery}

Approximately 40\% of global rabies cases occur in children under 15 years of age \citep{syn0232025}. Developing brains possess high structural and functional plasticity, characterized by elevated synaptic turnover, dendritic arborization, and cortical remapping potential \citep{james2026}. If viral clearance is achieved before irreversible structural neuronal loss occurs, pediatric neural circuits exhibit substantial capacity to reroute disrupted synaptic pathways.

\subsection{Epigenetic Mechanisms and BDNF Expression}

Recent pediatric neurotrauma research has elucidated dynamic epigenetic mechanisms of neurorecovery \citep{heinsberg2025}:
\begin{itemize}
    \item \textbf{BDNF DNA Hypomethylation}: Children recovering from acute brain injury exhibit significant DNA hypomethylation at the \textit{BDNF} promoter, enhancing transcription of Brain-Derived Neurotrophic Factor \citep{heinsberg2025}.
    \item \textbf{Rehabilitation-Induced Plasticity}: BDNF upregulation promotes long-term potentiation (LTP), dendritic remodeling, and synaptic repair, providing a validated target for early intensive neurorehabilitation \citep{james2026, heinsberg2025}.
\end{itemize}

\subsection{Neuroplastic Recovery Dynamics}

Longitudinal neuroimaging following pediatric neurosurgical interventions demonstrates that neural recovery follows a self-repair trajectory, progressing from subcortical structures to cortical networks over a 1-year window \citep{yan2025}. In the context of RVNA, early neurorehabilitation enhances adaptive neuroplastic changes and reduces long-term functional deficits \citep{james2026, yan2025}.

\section{Proposed Execution Workflow}

The RVNA Protocol is structured into six sequential clinical phases designed for execution in a neuro-intensive care environment:

\subsection{Phase 1: Diagnosis \& Triage (0--6 hours)}
\begin{enumerate}
    \item Confirm rabies diagnosis via real-time RT-PCR (saliva/nuchal skin biopsy) and direct fluorescent antibody (DFA) testing.
    \item Determine elapsed time from initial symptom onset and establish neurological staging.
    \item Obtain emergency compassionate-use regulatory and institutional ethical authorization \citep{fda2025}.
\end{enumerate}

\subsection{Phase 2: Catheter Placement (2--4 hours)}
\begin{enumerate}
    \item High-resolution 3T MRI stereotactic trajectory planning.
    \item Stereotactic placement of bilateral stepped-tip micro-cannulas into thalamic, basal ganglia, and pontine regions.
    \item Intraoperative CT/MRI verification of cannula trajectory.
\end{enumerate}

\subsection{Phase 3: Regional Hypothermia Induction (1--2 hours)}
\begin{enumerate}
    \item Initiate selective brain cooling to target $32\text{--}33\,^{\circ}\mathrm{C}$.
    \item Maintain systemic core temperature at $37\,^{\circ}\mathrm{C}$ using active warming blankets.
    \item Establish continuous intracranial pressure (ICP) and brain temperature telemetry.
\end{enumerate}

\subsection{Phase 4: Dual-Antibody Infusion (4--6 hours)}
\begin{enumerate}
    \item Prepare SYN023 ($0.3\,\mathrm{mg/kg}$ equivalent parenchymal volume) \citep{syn0232025}.
    \item Prepare anti-rabies biparatopic VHH cocktail ($50\text{--}100\,\mu\mathrm{g/kg}$) \citep{terryn2016}.
    \item Initiate co-infusion with $1.5\,\mathrm{mM}$ gadoteridol at $1.0\text{--}2.5\,\mu\mathrm{L/min}$.
    \item Monitor real-time parenchymal distribution ($V_d/V_i$) via intraoperative MRI \citep{narsinh2025, munjal2025}.
\end{enumerate}

\subsection{Phase 5: Monitoring \& Support (3--5 days)}
\begin{enumerate}
    \item Controlled rewarming at $0.25\,^{\circ}\mathrm{C/hr}$.
    \item Continuous neuro-ICU supportive care and EEG monitoring for subclinical seizures.
    \item Serial CSF viral load and biomarker quantification.
    \item Maintain cerebral perfusion pressure ($\mathrm{CPP} > 65\,\mathrm{mmHg}$).
\end{enumerate}

\subsection{Phase 6: Neurorehabilitation (Ongoing)}
\begin{enumerate}
    \item Initiate early multisensory and motor stimulation protocols.
    \item Longitudinal cognitive, speech, and physical rehabilitation \citep{james2026}.
    \item Monitor neuroplastic recovery and BDNF epigenetic profiles \citep{heinsberg2025}.
\end{enumerate}

\section{Ethical Justification and Regulatory Framework}

\subsection{Compassionate Use Framework}

Expanded access (compassionate use) provides an established regulatory pathway for patients with life-threatening conditions to receive investigational therapeutics outside of clinical trials when \textbf{no comparable or satisfactory alternative therapies exist} \citep{fda2025}.

The RVNA Protocol is formulated specifically for this population:
\begin{itemize}
    \item No established curative therapy exists for symptomatic rabies encephalomyelitis \citep{jackson2011}.
    \item Current standard of care is restricted to terminal palliative sedation \citep{lacy2024}.
    \item Untreated disease carries a virtually 100\% mortality rate.
\end{itemize}

\subsection{FDA Expanded Access Criteria}

Under FDA Expanded Access regulations (21 CFR Part 312.300 / Emergency IND pathways), experimental therapies may be authorized when the potential benefit justifies potential risks \citep{fda2025}:
\begin{quote}
``Investigational medical products have not yet been approved or cleared by FDA and FDA has not found these products to be safe and effective for their specific use. Furthermore, the investigational medical product may, or may not, be effective in the treatment of the condition, and use of the product may cause unexpected serious side effects.'' \citep{fda2025}
\end{quote}
In the absence of existing alternatives, the catastrophic natural history of rabies encephalomyelitis justifies consideration under expanded access.

\subsection{Pediatric Ethical Considerations}

Because approximately 40\% of global rabies deaths occur in children, the ethical imperative to provide active therapeutic options is paramount. Pediatric patients possess the greatest potential for neurodevelopmental functional recovery following acute insults; providing an active clearance pathway respects the patient's right to life-saving intervention over passive palliation.

\subsection{Informed Consent}

Informed consent protocols must explicitly delineate:
\begin{enumerate}
    \item The experimental nature and unproven clinical efficacy of the multi-modal stack;
    \item Known neurosurgical risks (intracranial hemorrhage, infection, edema);
    \item Alternative options (palliative comfort measures);
    \item The unconditional right of legal representatives to withdraw from the protocol at any stage; and
    \item Comprehensive data collection for international scientific learning.
\end{enumerate}

\section{Limitations \& Future Work}

\subsection{Limitations}

\begin{longtable}{@{}p{3.8cm}p{5.8cm}p{5.4cm}@{}}
\caption{Key limitations of the RVNA Protocol and proposed mitigations.}\label{tab:limitations}\\
\toprule
\textbf{Limitation} & \textbf{Description} & \textbf{Mitigation} \\
\midrule
\endfirsthead
\toprule
\textbf{Limitation} & \textbf{Description} & \textbf{Mitigation} \\
\midrule
\endhead
Untested Combination & The complete multi-modal protocol has not been evaluated in combined in vivo trials. & Prioritize staged rodent and non-human primate symptomatic challenge studies. \\
\addlinespace
Nanobody GMP Supply & Anti-rabies nanobodies (VHH) require dedicated GMP microbial expression batches. & Establish standardized \textit{Pichia pastoris} biomanufacturing pipelines \citep{terryn2016}. \\
\addlinespace
Neurosurgical CED Risks & Catheter insertion carries risks of intracranial hemorrhage, reflux, and infection. & Utilize stepped-design cannulas, micro-infusion rates ($<3.0\,\mu\mathrm{L/min}$), and real-time MRI guidance \citep{ced2025, narsinh2025}. \\
\addlinespace
Intracranial Pressure (ICP) & Parenchymal infusion into inflamed tissue risks elevating ICP and causing edema. & Continuous ICP telemetry, osmoprotective agents (mannitol / hypertonic saline), and ramped flow protocols. \\
\addlinespace
Animal Model Gap & Lack of established large-animal models replicating symptomatic CNS viral clearance. & Implement systematic feline or canine translational models alongside compassionate clinical registry data. \\
\addlinespace
Viral Load Quantification & Real-time measurement of parenchymal viral load and synaptic saturation is challenging. & Develop serial CSF biomarker panels and microdialysis sampling protocols. \\
\addlinespace
Thermal Gradient Control & Selective brain cooling requires precise temperature titration without systemic shivering. & Closed-loop nasopharyngeal / endovascular heat exchangers with active core warming blankets ($37\,^{\circ}\mathrm{C}$). \\
\addlinespace
Post-Lysis Inflammation & Rapid extracellular virion neutralization may trigger local neuroinflammatory cascades. & Adjunctive neuroprotective modulation and controlled hypothermic dampening of cytokine release. \\
\bottomrule
\end{longtable}

\subsection{Future Work}

\begin{enumerate}
    \item \textbf{Preclinical In Vivo Validation}: Evaluate combined RVNA components (CED + regional hypothermia + IgG + biparatopic VHH) in symptomatic animal challenge models.
    \item \textbf{Nanobody Optimization}: Engineer multivalent and biparatopic nanobodies targeting conserved viral glycoprotein epitopes \citep{terryn2016}.
    \item \textbf{Computational Fluid Dynamics (CFD) Modeling}: Model CED distribution dynamics in infected cerebral tissue to optimize stereotactic trajectory planning \citep{ced2025}.
    \item \textbf{Compassionate Use Clinical Registry}: Establish an international multi-center registry to document outcomes of individual compassionate-use cases \citep{fda2025}.
    \item \textbf{Formulation Optimization}: Develop specialized nanocarrier and hydrogel formulations for stabilized intracranial CED antibody delivery \citep{ced2025}.
\end{enumerate}

\section{Conclusion}

\begin{quote}
``Although preventative therapy for rabies is highly successful after recognized exposures, there is no established therapy that is effective for patients who develop rabies encephalomyelitis, which remains an important therapeutic challenge for the future.'' \citep{jackson2011}
\end{quote}

For decades, the clinical management of symptomatic rabies has remained trapped between therapeutic despair and unscientific coma protocols. The RVNA Protocol proposes a fundamental paradigm shift: replacing passive metabolic sedation with a targeted, biophysically grounded, active clearance architecture.

By combining Convection-Enhanced Delivery, dual full-IgG and synaptic-penetrating nanobodies, targeted regional brain hypothermia, and pediatric neuroplastic potential, the RVNA framework provides an actionable blueprint for translational investigation and compassionate clinical trial in medicine's most lethal infectious challenge.

\begin{quote}
``More basic research is needed to improve our understanding of the mechanisms involved in rabies pathogenesis, which will hopefully facilitate the development of novel therapeutic approaches in the future.'' \citep{jackson2011}
\end{quote}

The RVNA Protocol is a contribution to that future.

\section*{Acknowledgements}
The authors thank the global community of neurovirologists, neurosurgeons, and biomedical engineers whose open scientific contributions in rabies pathogenesis, convection-enhanced delivery, antibody engineering, and therapeutic hypothermia informed this systems architecture. No external funding was received for this work.

\section*{Medical and Legal Disclaimer}
This document describes a theoretical biomedical systems engineering proposal for scientific research and expanded-access evaluation. It does not constitute formal medical advice. All clinical procedures, expanded-access applications, and surgical interventions must be evaluated and administered by licensed medical specialists in accordance with applicable national laws, institutional review boards (IRBs), and regulatory health authorities.

\bibliography{citations}

\end{document}
